% P57 | 7 Conclusions | Limitations and next steps
% guide: dev_notes/paragraph_guides/P57_limitations_and_next_steps.txt
The model assumes a constant volatility, whereas real markets show an implied-volatility smile and term structure, so a local- or stochastic-volatility model would be the natural next step.
The tree would then no longer recombine so simply, and the PDE would gain a second dimension.
The inputs are simplified: a historical volatility instead of an implied one, a Canadian rate for a US-dollar stock, and no dividends for NVIDIA.
In practice each of these would be replaced by market data, and the sensitivities of Section~\ref{sec:B-greeks} say how much the price would move.
There are no transaction costs and no jumps in the model, and discrete hedging is studied only for the call.
Jump or gap risk would leave a residual that no rebalancing frequency removes.
Natural next steps are to calibrate to an option chain, to add a discrete dividend schedule, and to extend the regression Monte Carlo to a basket of underlyings.
